
\lussec 9. $2+1$ 味 QCD 中的示例计算

以下几段的目的是演示所提出的
计算重正化常数 $\ZA$、改善系数 $\cfl$
以及低能常数 $\Sigma,F$ 诸策略的可行性。
这一检验只需一次模拟运行即可完成，
并且并不试图真正计算这些低能常数，
因为那将需要把数据外推到无限体积、手征与连续极限。


\lussubsec 9.1 模拟细节

下面报告的计算基于最近一次对 $2+1$ 味夸克 QCD 的模拟
[\cite{TMopenQCD}]。
在该运行（文献~[\cite{TMopenQCD}] 中标记为 $I_1$）中，
模拟了一个在时间方向具有开边界条件 [\cite{openQCD}]
的 $64\times32^3$ 格点，
其参数空间中的位置对应于
估计为 $0.090$ fm 的格距
[\cite{AokiEtAlI},\cite{AokiEtAlII}]，
$\pi$ 介子与 K 介子质量分别约为 203 与 520 MeV
[\cite{TMopenQCD}]。
因此格点的物理尺寸使得
所计算的介子质量与矩阵元的有限体积效应预期很小。

对于该模拟使用的 Iwasaki 规范作用量 [\cite{Iwasaki}]
及标准 O($a$) 改进 Wilson--Dirac 算符
[\cite{SW},\cite{SFimp}]，
系数 $\csw,\ca$ 已非微扰地已知
[\cite{AokiEtAlIII},\cite{KanekoEtAl}]。
该运行产生了由 150 个规范场组态构成的代表性系综，
相邻组态相隔 $9.6$ 个分子动力学时间单位。
为了压制残余的自关联效应，
本节所引的统计误差通过如下方式估计：
把基本观测量系列的测量数据分成每 5 个连续测量为一组的 bin，
再对分箱数据应用 jackknife 方法。

% Fig. 1 (converted from epsf/figure1; graphic from plots/wi.eps -> images/wi.pdf)
\topinsert
\includegraphics[width=10.5cm]{images/wi}
\vspace{0.3cm}
\figurecaption{%
PCAC 关系 (8.8) 中出现的关联函数组合，
绘成求和范围 $d$ 的函数。
所示关联函数的线性组合分别是乘以重正化常数 $\ZAh{ud}$ 的（三角）、
乘以改善系数 $\cfl$ 的（菱形）
以及乘以因子 $1-\ZAh{ud}\ctp m_{ud}$ 的（方块）那些组合，
流时间为 $t=3.86$。
当 $d\geq9$ 时，函数 $\CAtP{ud}$（圆）比其他函数小许多个数量级，
可以安全地忽略（参见 8.1 小节）。
所有量均以格点单位给出。
}
\endinsert


\lussubsec 9.2 $\ZAh{rs}$ 与 $\cfl$ 的计算

进入 PCAC 关系 (8.7) 的关联函数
在流时间 $t$（以格点单位计）取值
$2.47,\,3.86,\,5.56,\,7.56$ 处被计算，
对应的梯度流光滑范围分别为 $0.4,\,0.5,\,0.6,\,0.7$ fm。
为了避免靠近格点在时间方向的边界，
随时间演化场 $P_t^{rs}(y)$ 被插在格点中部。
随后利用时刻 $y_0$ 处的 12 个随机源场，
对各关联函数作空间平移平均。

PCAC 关系中的全部关联函数
都能在整个求和窗口 $d$ 的范围内以很小的统计误差得到。
正如所预期的，
当 $d$ 大于梯度流光滑范围的 2 或 3 倍时，
乘以轴矢流重正化常数的各关联函数之和
与其他关联函数都达到平台
（代表性情形见图~1；
乘以 $\ZAh{ud}$ 的函数在接近格点边界处的轻微下弯
是一种格距效应）。

按照第 8 节概述的策略，
现在通过要求 PCAC 关系 (8.8)
在 $t$ 的两个值 $t_1,t_2$ 处精确成立来确定比值 (8.9)。
只要处于各关联函数的平台范围内，
求和窗口 $d$ 的选取对算得的比值影响很小。
取 $d=18$ 便得到表~1 第二、三列所引的结果，
其中因子 $1-\ZAh{ud}\ctp m_{ud}$
用系数 $\ctp$ 的树图值 (6.22) 作了估计。
不出所料，这一修正非常小（约 $0.2\%$）。

在上-奇异味道道中有两个 O($a$) 质量修正来源。
照上述步骤处理，
来自因子 $1-\ZAh{us}\ctp m_{us}$ 的修正估计约为 $1.3\%$，
而来自式 (8.10) 中正比于 $b_{\chi}$ 项的修正
使结果改变约 $0.2\%$。
如表~1 所示，
两个味道道所得 $\cfl$ 的值之间有显著的一致性。
对流时间的依赖也很小，
而重正化常数 $\ZAh{ud}$ 与 $\ZAh{us}$ 的差别
很可能完全可以由式 (8.4) 中的依赖质量的因子解释。
O($a$) 改进于是进行得非常好，
就重正化常数而言残余格距效应仅为百分之一的几分之一水平。

% Table 1 (re-typeset as LaTeX tabular; data verbatim from source table1)
\topinsert
	ablecaption{$\ZAh{rs}$ 与 $\cfl$ 的模拟结果}
\begin{center}
\renewcommand{\arraystretch}{1.25}
\begin{tabular}{ccccc}
\hline\hline
$t_1,t_2$ & $\ZAh{ud}$ & $\cfl$ & $\ZAh{us}$ & $\cfl$ \\
\hline
$2.47,3.86$ & $0.7929(87)$ & $0.5703(11)$  & $0.8096(41)$ & $0.56867(60)$ \\
$2.47,5.56$ & $0.7904(85)$ & $0.5712(14)$  & $0.8073(40)$ & $0.56990(72)$ \\
$2.47,7.56$ & $0.7895(85)$ & $0.5716(16)$  & $0.8067(39)$ & $0.57025(83)$ \\
$3.86,5.56$ & $0.7885(85)$ & $0.5734(22)$  & $0.8054(40)$ & $0.5726(10)$  \\
$3.86,7.56$ & $0.7883(85)$ & $0.5735(26)$  & $0.8055(39)$ & $0.5724(12)$  \\
$5.56,7.56$ & $0.7882(86)$ & $0.5738(36)$  & $0.8056(38)$ & $0.5723(16)$  \\
\hline\hline
\multicolumn{5}{l}{\footnotesize The values of $\cfl$ in the 3rd and 5th column}\\[-3pt]
\multicolumn{5}{l}{\footnotesize were obtained in the $ud$ and $us$ channel, respectively}\\
\end{tabular}
\end{center}
\endinsert


表~1 所引的重正化常数 $\ZAh{ud}$ 的结果
与文献~[\cite{AokiEtAlIV}] 早先使用基于薛定谔泛函的方法
[\cite{ZASFI},\cite{ZASFII}]
得到的值 $\ZA=0.781(20)$ 一致。
注意改善系数 $\cfl$ 结果上接近其树图值 (6.22)，
尽管在所考虑的规范耦合处并没有很好的理由必然如此。


\lussubsec 9.3 手征凝聚的计算

确定了改善系数 $\cfl$ 之后，
未重正化的随时间演化凝聚 $\Sigma^{uu}_t$
便可按 7.2 小节所述直截了当地计算。
源时刻 $x_0$ 取在格点中部，
以尽量减小激发态对标量密度期望值的贡献。
使用 12 个随机源场并在 $t_1,t_2=3.86,5.56$ 处取 $\cfl$，
计算给出表~2 第二列所引的结果。
在所考虑的流时间范围内，
可见随时间演化的凝聚随 $t$ 增大而减小，粗略如 $1/t$。
统计误差跟随这一演变，
无论如何都小得令人鼓舞。

% Table 2 (re-typeset as LaTeX tabular; data verbatim from source table2)
\topinsert
	ablecaption{格点单位下 $\Sigma^{uu}_t$ 与 $G^{ud}_t$ 的模拟结果}
\begin{center}
\renewcommand{\arraystretch}{1.25}
\begin{tabular}{ccccc}
\hline\hline
$t$ & $\Sigma^{uu}_t$ & $G^{ud}_t$
    & $\Sigma^{uu}_tG^{ud}/G^{ud}_t$ & $\Sigma^{uu}_t/G^{ud}_t$ \\
\hline
$2.47$ & $0.0006277(95)$ & $0.01484(23)$  & $0.003962(61)$ & $0.04230(85)$ \\
$3.86$ & $0.0004251(58)$ & $0.01028(16)$  & $0.003872(55)$ & $0.04134(81)$ \\
$5.56$ & $0.0002911(36)$ & $0.00720(11)$  & $0.003785(51)$ & $0.04040(78)$ \\
$7.56$ & $0.0002017(23)$ & $0.005092(80)$ & $0.003711(48)$ & $0.03962(76)$ \\
\hline\hline
\end{tabular}
\end{center}
\endinsert


为了能把随时间演化的凝聚与手征极限下的
夸克凝聚 $\Sigma$ 联系起来，
还需要计算真空到 $\pi$ 介子矩阵元
\equation{
  \Gpi=\ZP(1+b_P\mq{ud}+\bar{b}_P\tr\Mq)\kern1pt G^{ud},
  \enum
  \nexteq{3.0ex}
  \Gpit=Z_{\chi}(1+b_{\chi}\mq{ud}+\bar{b}_{\chi}\tr\Mq)\kern1pt G^{ud}_t,
  \enum
}
实际上，由于式 (4.24)、(4.25) 中只出现比值
\equation{
 {\rens{\Sigma}{t}\over\Gpit}=
 {\Sigma^{uu}_t\over G^{ud}_t}
 \enum
}
故只需计算 $\Gpi$ 及未重正化的矩阵元 $G^{ud}_t$。
后者可以在大时间间隔 $|x_0-y_0|$ 处
从赝标关联函数 (7.1) 确定，
其方式与通常从普通赝标两点函数提取流时间为零处的
裸矩阵元 $G^{ud}$ 相同。
把 $x_0$ 取在紧邻格点边界的值处，
在所有情形下都存在一个较宽的 $y_0$ 时间范围，
其中单 $\pi$ 介子中间态占主导，
因而可以容易地确定所需的矩阵元。

与随时间演化的凝聚类似，
矩阵元 $G^{ud}_t$ 是 $t$ 的单调递减函数。
然而表~2 最后两列所列的比值随 $t$ 的变化小得多，
这与如下事实一致：
在手征极限下这些比值对流时间的依赖必须消失。


\lussubsec 9.4 若干说明

\noindent{\itshape 换算到物理单位.}
模拟格点的格距 $a=0.08995(40)$ fm 由 PACS-CS 通过计算
$\Omega$ 重子质量确定 [\cite{AokiEtAlII}]。
PACS-CS 还利用步长标度与薛定谔泛函计算了把裸矩阵元 $G^{ud}$
与 $\MSbar$ 方案下 2 GeV 处的 $\Gpi$ 相联系的重正化常数
$\ZP=0.580(21)$ [\cite{AokiEtAlIV}]\kern1pt\sfn{$\dagger$}{%
此处给出的 $\ZP$ 值需要把文献~[\cite{AokiEtAlIV}]
所引的若干表项组合起来并将误差按平方和合成而得到。}。
例如在 $t=3.86$ 处即得结果
\equation{
  {\rens{\Sigma}{t}\Gpi\over\Gpit}=[287(4)\,\MeV]^3,
  \enum
  \nexteq{2.5ex}
  {\rens{\Sigma}{t}\over\Gpit}=91(2)\,\MeV,
  \enum
}
其中忽略了式 (9.1) 中的 O($a$) 质量修正
（反正式 (9.4) 的误差由 $\ZP$ 的误差主导）。
显然，虽然这些比值可能已接近手征极限下的凝聚 $\Sigma$
与衰变常数 $F$，
但要在可控系统误差下确定 $\Sigma$ 与 $F$，
还需要对一系列格距进行模拟。

\vskip0.5ex
\noindent{\itshape 赝标衰变常数.}
由于轴矢流重正化常数的计算不需要单独的模拟，
重正化的赝标介子衰变常数在每个被模拟的格点上都是直接可得的。
不过根据味道道以及所要达到的精度，
用于确定重正化常数 $\ZAh{rs}$ 的 PCAC 关系中的
依赖质量 O($a$) 修正可能需要非微扰地估计。

\vskip0.5ex
\noindent{\itshape 向连续极限的标度.}
本节报告的计算结果依赖于流时间的选择。
当模拟不同格距的格点时，
应当把这些参数以某个合适的参考标度为单位固定，
例如参考流时间 $t_0$ [\cite{WilsonFlow}]。
只有这样，才能期望所计算的重正化量
以正比于 $a^2$ 的速率收敛到连续极限。
